\subsection{Reutilización de Código}

\subsubsection{Regla: principio DRY}

\textbf{Regla:} Se debe cumplir estrictamente el principio DRY (\emph{Don't Repeat Yourself}). Queda prohibido duplicar un bloque de código: si una misma lógica se repite en dos o más métodos, deberá extraerse a un método privado reutilizable.

\textbf{Descripción:} Hunt y Thomas formularon el principio DRY como la exigencia de que todo conocimiento tenga una representación única, inequívoca y autorizada dentro de un sistema, de manera que un cambio de regla solo deba aplicarse en un lugar (Hunt \& Thomas, 2019). Fowler describe la extracción de un fragmento duplicado hacia un método con nombre propio como una de las refactorizaciones más elementales para eliminar duplicación (Fowler, 2018).

\textbf{Código correcto:}
\begin{lstlisting}[style=csharpstyle]
public sealed class MoveValidator
{
    public bool CanMoveTo(Player player, Tile destination)
    {
        return IsWithinBoardBounds(destination) && !destination.IsOccupied;
    }

    public bool CanTeleportTo(Player player, Tile destination)
    {
        return IsWithinBoardBounds(destination) && !destination.IsOccupied;
    }

    private bool IsWithinBoardBounds(Tile tile)
    {
        return tile.Row >= 0 && tile.Column >= 0;
    }
}
\end{lstlisting}

\textbf{Código incorrecto:}
\begin{lstlisting}[style=csharpstyle]
public sealed class MoveValidator
{
    public bool CanMoveTo(Player player, Tile destination)
    {
        return destination.Row >= 0 && destination.Column >= 0 && !destination.IsOccupied;
    }

    public bool CanTeleportTo(Player player, Tile destination)
    {
        return destination.Row >= 0 && destination.Column >= 0 && !destination.IsOccupied;
    }
}
\end{lstlisting}
